\documentclass[10pt,a4paper]{amsart}
\usepackage[margin=19mm]{geometry}
\usepackage{amsmath,amssymb,mathtools}
\usepackage[scr=rsfs,cal=euler]{mathalfa}
\usepackage{xfrac}
\usepackage[unicode,hidelinks,bookmarks=false]{hyperref}
\newcommand{\ie}{i.e.\ }
\newcommand{\cf}{cf.\ }
\newcommand{\resp}{respectively\ }
\newcommand{\Subsection}{\subsection}
\providecommand{\nobreakdash}{\nobreak\hbox{-}}
\setlength{\emergencystretch}{2em}
\multlinegap0pt
\usepackage{bm}
\usepackage{bbm}
\usepackage{dsfont}
\usepackage{xspace}
\usepackage{cancel}
\usepackage{mathtools}
\usepackage{amsthm}
\makeatletter
\@ifundefined{theorem}{\newtheorem{theorem}{Theorem}}{}
\@ifundefined{proposition}{\newtheorem{proposition}[theorem]{Proposition}}{}
\makeatother
\title[Second-Order Optimality Conditions for Nonsmooth Constrained]{Second-Order Optimality Conditions for Nonsmooth Constrained Optimization with Applications to Bilevel Programming}
\author{Xiang Liu, Mengwei Xu, Liwei Zhang}
\date{Source: arXiv:2511.02439v1 (CC BY 4.0)}
\begin{document}
\maketitle

\noindent\emph{Source and licence.} Second-Order Optimality Conditions for Nonsmooth Constrained Optimization with Applications to Bilevel Programming. Version arXiv:2511.02439v1 (CC BY 4.0), distributed under the Creative Commons Attribution 4.0 International licence, as recorded in the arXiv licence metadata for this exact version and shown on the abstract page https://arxiv.org/abs/2511.02439v1. Full rights notice: licenses/arxiv-2511.02439-CC-BY-4.0.txt.\par\smallskip

\noindent\textit{Editorial scope.} Counted unit: the Theorem whose source label is \texttt{thm SOSC for
nonsmooth} in the article (source0.tex lines 975--991, source label \texttt{thm SOSC for
nonsmooth}), delivered together with the article's own complete original proof of it (source0.tex lines 992--1043, one \texttt{proof} environment of 52 lines). No mathematics is written, repaired or re-derived: statement and proof are the article's own text line for line. Required context retained uncounted: prop lipschitz--Hadamard directional differentiable (lines 319--325). Every definition, display and label of those lines is retained unchanged. Omitted from this package and not redistributed: 1--318, 326--974, 1044--2146. Nothing inside a retained excerpt is omitted or altered: every retained body is delivered exactly as the article's lines are written, so no omission receipt of any kind accompanies this package. Citations: the retained excerpts carry no citation command at all, so no bibliography entry of the article is transcribed and no reference is redistributed. Reference adaptation: none was needed and none was made; every reference inside the delivered unit resolves inside it, so no unresolved reference or citation is produced. Printed slips kept literal: no printed word, symbol, hypothesis or number of the counted statement or of its proof was altered. Layout and class: the article's own class is \texttt{article} with packages including graphicx, amssymb, amsbsy, ifthen, cite, amsfonts, amsmath, amsthm, geometry, algorithm2e, xcolor, verbatim, dsfont, hyperref; the wrapper is the lane's pinned \texttt{amsart} editorial wrapper at 10pt on A4, which restates the theorem-like environments the retained text needs and adds the article's own macro definitions that the retained text needs (none), with extra wrapper packages (bm, bbm, dsfont, xspace, cancel, mathtools). The delivered numbering is the unit's own. Assets: no retained span contains a figure environment, a graphics-inclusion command, a PDF-page inclusion, a table or a TikZ picture; no image file is copied into this package, the package declares no asset dependency and no image file is redistributed with it. Licence: version arXiv:2511.02439v1 of this article is distributed under the Creative Commons Attribution 4.0 International licence, as recorded in the arXiv licence metadata for this exact version and shown on the abstract page https://arxiv.org/abs/2511.02439v1. A full rights notice is delivered at licenses/arxiv-2511.02439-CC-BY-4.0.txt. Characters: every retained excerpt is pure ASCII, so the delivered engine, XeLaTeX with its default Latin Modern text font, reports no missing character.
\begin{proposition}\label{prop lipschitz--Hadamard directional differentiable}
  Suppose that $g:\mathbb{R}^n\to \mathbb{R}^m$ is second-order directionally
differentiable at $x^*$ in direction $d$ and locally Lipschitz continuous
 with  modulus $c$ at $x^*\in\mathbb{R}^n$, then $g$ is second-order directionally differentiable at $x^*$ in
direction $d$, in the Hadamard sense and
the second-order directional derivative $g''(x^*;d,\cdot)$ is Lipschitz continuous with modulus $c$ on $\mathbb{R}^n$.
\end{proposition}

\begin{theorem}[Second-order sufficient optimality conditions]\label{thm SOSC for
nonsmooth}
Let $x^*$ be a feasible point of $(P)$. Suppose that $f$ is  %locally Lipschitz continuous, second-order directionally differentiable and
    second-order
 epi-regular at $x^*$ and the set $\Psi$ is outer second-order regular at $x^*$ in every direction  $d\in T_{\Psi}(x^*)$.
 Assume that  we have
 \begin{align}\label{thm f' ge 0}
   & f'(x^*;d)\ge 0, \quad \forall d\in T_{\Psi}(x^*),
 \end{align}
 and  the optimization problem
   \begin{align}\label{thm p f''}
    \min_{w}~ f''(x^*;d,w) ~~~{\rm s.t.}~w\in \mathcal{T}_{\Psi}^2(x^*,d),
   \end{align}
possesses a positive optimal objective value for every $d\in C_{\Psi}(x^*):=\left\{d\in T_{\Psi}(x^*): f'(x^*;d)\le
   0\right\}  \backslash \{0\}$. Then the
second-order growth condition holds for $(P)$ at $x^*$.
\end{theorem}

\begin{proof}
 Assume to the contrary that the second-order growth condition does not hold for $(P)$ at $x^*$.
  Then there exists a sequence $x^k\in \Psi$, $x^k\neq x^*$, converging to $x^*$ such that
\begin{align}\label{f contradiction}
  & f(x^k)-f(x^*)\le o(\|x^k-x^*\|^2).
\end{align}
  Select a subsequence if necessary, we define $t_k=\|x^k-x^*\|$, $t_k\downarrow 0$ and $t_k>0$. Let
$d^k=\frac{x^k-x^*}{t_k}$, then $x^k=x^*+t_k d^k$ and $\|d^k\|=1$. Without loss of generality, we have that $d^k\to d\in T_{\Psi}(x^*) $
and $\|d\|=1$.

Since $f$ is locally Lipschitz continuous and directionally differentiable at $x^*$,
then
\begin{align*}
   & f'(x^*;d)=\lim_{k\to \infty}\frac{f(x^*+t_k d^k)-f(x^*)}{t_k}
   \le \lim_{k\to \infty}\frac{o(\|x^k-x^*\|^2)}{t_k}=0,
\end{align*}
thus $d\in \mathcal{C}_{\Psi}(x^*)$ and by (\ref{thm f' ge 0}), we have $f'(x^*;d)=0$.
Let $w^k=\frac{x^k-x^*-t_k d}{\frac{1}{2}t_k^2}$, then $x^k=x^*+t_k d+\frac{1}{2}t_k^2w^k\in \Psi$ and $t_kw^k\to
0$. Since $\Psi$ is outer second-order regular
at $x^*$ in  direction $d\in T_{\Psi}(x^*)$, we have
 \begin{align}\label{d(w^k,T^2) to 0}
   &d(w^k,\mathcal{T}^2_{\Psi}(x^*,d))\to 0.
 \end{align}
  Since $f$ is second-order  epi-regular at $x^*$, then
 $$f(x^*+t_k d+\frac{1}{2}t_k^2w^k) \ge f(x^*)+t_k f'(x^*;d)+\frac{1}{2}t_k^2 f''(x^*;d,w^k)+o(t_k^2).$$
 By
 (\ref{f contradiction}) and the fact $f'(x^*;d)=0$, we have
 \begin{align}\label{limsup f''(x^*;d,w^k)}
   & \limsup_{k\to \infty}f''(x^*;d,w^k)\le 0.
 \end{align}
 Notice that the optimal value of Problem (\ref{thm p f''}) is positive,
 there exists some constant $\beta>0$ such that
 \begin{align*}
    & \forall w\in \mathcal{T}^2_{\Psi}(x^*,d),~~~f''(x^*;d,w)\ge \beta >0.
 \end{align*}
It follows from Proposition \ref{prop lipschitz--Hadamard directional differentiable} that $f''(x^*;d,\cdot)$ is
Lipschitz continuous, then there exists some constant $L_f>0$ such that
\begin{align}\label{Lf}
   & |f''(x^*;d,w^k)-f''(x^*;d,w)|\le L_{f}\|w^k-w\|.
\end{align}
By (\ref{d(w^k,T^2) to 0}), there exist $v^k\in \mathcal{T}^2_{\Psi}(x^*,d)$ such that $\|w^k-v^k\|\le
\frac{\beta}{2L_f}$ for all $k$ large enough.
Formula (\ref{Lf}) yields
\begin{align*}
   & f''(x^*;d,w^k)\ge f''(x^*;d,v^k)-L_f\|w^k-v^k\|\ge \beta-\frac{\beta}{2}=\frac{\beta}{2}>0,
\end{align*}
%thus we obtain
%\begin{align*}
%   & \liminf_{n\to \infty}f''(x^*;d,w^k)\ge \frac{\beta}{2}>0.
%\end{align*}
 which contradicts with (\ref{limsup f''(x^*;d,w^k)}). We complete the proof.
\end{proof}

\end{document}
